\section{Further questions}\label{sec:questions}

The theorem determines the maximizing body and the angular reach of
its passages. It leaves several questions about the proof and nearby
geometric problems.

\runin{A quantitative equality theorem.}
If an ambidextrous sofa has area at least $\Mc-\eps$, how close must
it be, after a rigid motion, to $\SR$? The exact partition deficits
and the coercive support-function remainder are possible starting
points. A Hausdorff estimate would also have to control the geometry
of nearly maximizing caps before they satisfy the exact maximizer
source laws. We do not claim such an estimate here.

\runin{The ordinary-sofa input.}
Two strict exclusions use the ordinary-sofa bound $22199/10000$
from Baek's theorem. A direct ambidextrous estimate strong enough to
replace those uses could shorten the dependency of the proof.
Removing that input requires estimates for the same genuinely
feasible auxiliary bodies; replacing them by arbitrary envelopes
would not justify the ordinary-sofa theorem.

\runin{Passage paths.}
The quarter-turn theorem does not classify all motions of the
optimal body. A classification would need to account for waiting,
time reparametrizations, translations in straight arms, and any
additional angular excursions. The explicit paths in
\cref{eq:reference-motion} prove existence but do not provide that
classification.

\runin{Other corridors.}
The common incoming orientation, equal widths, and right angles are
part of the problem proved here. Unequal hallway widths and other
corner angles change both the supporting-wall formulas and the
terminal constraints. The cap method supplies a framework for those
questions, but the sharp calibration and equality shape would need
to be determined again.

For a hallway width $d>0$, scaling the present definitions gives
the immediate consequence that the largest ambidextrous area is
$d^2\Mc$, attained exactly by the scaled rigid images of $\SR$.
This follows by applying $p\mapsto p/d$ to both the body and the
passages.

\runin{Variation methods with moving contact sets.}
The polygon approximation, curvature estimates, and graph-area
identities used here may apply to other extremal problems in which
contact regions change with the shape. A general formulation would
need to retain the moving integration boundaries and the shared
occupation of the two inner walls at a corner.
